%!TEX root = ../main.tex

\section{Front End}
According to the use cases a user should be able to monitor data created on the network.
This section describes the test done to verify that this functionality is implemented correctly.
\subsection{Method}
Figure \ref{fig:frontendsetup} depicts an overview of the setup used in the verification.
An amount of GPS data is recorded and is presented to the GPS node using service virtualisation.\\

The two parts, CAN-bus and WiFi node were not finished, see sections \ref{sec:methods_to_implement_can} and \ref{sec:somewifinodesection} respectively, as such it was necessary to create service virtualisation for this link.
A small utility was written to serve this purpose, it virtualises the dotted boxes seen on figure \ref{fig:frontendsetup}.
This utility reads the messages sent by the GPS node, extracts the timestamp and inserts it into the data frame in place in the DLC nibble and then outputs it to stdio.
The message is then on the form seen in figure \ref{fig:backendmsg}, the form expected by the front end.
This means that, seen from the perspective of the front end the underlying structure is correct.

\begin{figure}
	\centering
	\includegraphics[width=.75\linewidth]{graphics/frontend_verification}
	\caption{The setup used to verify the functionality of the front end. 
	The dotted boxes are virtualised but the correct connection is maintained between the Zybo and the PC.}
	\label{fig:frontendsetup}
\end{figure}

In order to most accurately reproduce the actual function of the system, the test was done using a wireless connection between a Zybo and a PC.
The connection was established using the method described in section \ref{sec:wifi}.
On the PC, the front end was started, taking its input from a socat listener.
\begin{lstlisting}
>> ./frontend | socat - tcp-listen:2049
\end{lstlisting}
Then, on the Zybo, the GPS and WiFi nodes were started, piping their output to socat:
\begin{lstlisting}
>> ./sensornode | ./verification | socat - tcp:196.178.10.10:2049
\end{lstlisting}
Here, the first and last data points from the data set.
\begin{lstlisting}
>> N 55.36734 E 10.43185
>> N 55.36707 E 10.43109
\end{lstlisting}
This corresponds with the data created on the GPS node.
